\section{GTM 7: a course in arithmetic}

Obey to the tradition of this book by Serre, we will assume that the characteristic of \(k\) is not two.

\subsection{General quadratic forms}

\begin{definition}
    Let \(V\) be a module over a commutative ring \(R\), a quadratic form is a map \(q: V \to R\)  such that
    \begin{enumerate}
        \item \(q(av) = a^2q(v)\) for any \(a \in R\) and \(v \in V\).
        \item \(B_q(v,w) = q(v+w) - q(v) - q(w)\) is a bilinear form on \(V\).
    \end{enumerate}
    and we call the pair \((V,q)\) to be a quadratic module over \(R\).
\end{definition}

we set the scalar product to be 
\[x \cdot y = \frac{1}{2}B_q(x,y) = \frac{1}{2}(q(x+y) - q(x) - q(y))\]
for any \(x,y \in V\), such a definition is similar to the inner product of a vector space, it makes sense if the ring is not of characteristic two. Moreover, we restrict our discussion of module to the finitely generated free \(R\)-modules, it will be nice to talk about the matrix or derminant.


\begin{remark}
    By definition
    \[x \cdot x = \frac{1}{2}(q(2x)-q(x)-q(x)) = \frac{1}{2}(4q(2x)-2q(x)) = q(x)\]
    and if \(\{e_1,...,e_n\}\) is a basis of \(V\), then we can write 
    \[A = (e_i \cdot e_j)_{i,j}\]
    to denote the matrix of quadratic form \(q\) with respect to the basis, then we have
    \[q(x) = x^tAx\]
    such a matrix is symmetric and gives a one-to-one correspondence
    \[\text{quadratic forms} \longleftrightarrow \text{symmetric bilinear forms}\]
\end{remark}

\begin{remark}
    We use \(\mathbf{Quad}_R\) to denote the category of quadratic forms over \(R\), the objects are quadratic modules \((V,q)\); If \((V,q)\) and \((V',q')\) are two objects, then the (metric) morphism \(f:V \to V'\) is a \(R\)-linear map such that \(q'(f(v)) = q(v)\) for any \(v \in V\), i.e. \(f\) preserves the quadratic form.
    % https://q.uiver.app/#q=WzAsMyxbMCwwLCJWIl0sWzEsMCwiViciXSxbMCwxLCJSIl0sWzAsMiwicSIsMl0sWzEsMiwicSciXSxbMCwxLCJmIl1d
\[\begin{tikzcd}
	V & {V'} \\
	R
	\arrow["f", from=1-1, to=1-2]
	\arrow["q"', from=1-1, to=2-1]
	\arrow["{q'}", from=1-2, to=2-1]
\end{tikzcd}\]

a remarkable property (isometry) is that for any \(x,y \in V\), we have
\[f(x) \cdot f(y) = x \cdot y\]
\end{remark}

\begin{remark}
    If we change the basis \(\{e_1,\dots,e_n\}\) by a invertible matrix \(P\), then the matrix of quadratic form under the new baisis \(\{Pe_1,\dots,Pe_n\}\) is given by
    \[A' = P^T A P\]
    we define the \textbf{discriminant} of quadratic form \((V,q)\) to be the determinant of the matrix \(A\), so by above discussion, the discriminant is unique up to multiplication by an element of \(R^2\) since \(\det A' = \det A \cdot \det P^2\).
\end{remark}

\begin{remark}
    The explicit form of quadratic form is of the form
    \[q(X_1,\dots,X_n) = \sum_i c_{i,i}X_i^2 + 2\sum_{i<j}c_{i,j}X_iX_j\]
    or just of the form 
    \[q(X_1,\dots,X_n) = \sum_{i,j} a_{i,j} X_i X_j\]
    In the first form, the martix of quadratic form is just \(A = (c_{i,j})\), and sure it is symmetric.

    Two quadratic forms \((V,q)\) and \((V',q')\) are called \textbf{equivalent} if \(V\) isomorphic to \(V'\) and there exists a metric morphism \(T: V \to V' \) such that \(q' \circ T = q\), or equivalently, the discriminants of \(q\) and \(q'\) are equal modulo \(R^2\) and the rank of \(q\) and \(q'\) are equal.
\end{remark}

\textcolor{blue}{Now we talk about the orthogonality of quadratic forms.}

\begin{definition}
    Let \((V,q)\) be a quadratic module over \(R\).

    \begin{itemize}
        \item If \(x,y \in V\) and \(x \cdot y = 0\), then we say that \(x\) and \(y\) are \textbf{orthogonal}.
        \item If \(H \subset V\) is a subset, and we define 
        \[H^\perp : = \{x \in V \mid x \cdot y =0 \text{ for } \forall y \in H\}\]
        such a subset is always a submodule of \(V\).
        \item If \(V_1\) and \(V_2\) are two submodules of \(V\), they are said to be \textbf{orthogonal} if \(x \cdot y =0 \) for any \(x \in V_1\) and \(y \in V_2\), or equivalently, \(V_1 \subset V_2^\perp\) .
        \item \(V^\perp \) is called the \textbf{radical} of \(V\), denoted by \(\mathrm{rad}(V)\), and the rank of the quadratic form is defined as the codimension of the radical space:
        \[\mathrm{rank}\, q : = \dim V - \dim V^\perp\]
    \end{itemize}
\end{definition}

suppoese that \(R = k\) is a field, in the context of vector space we can consider dual space, and if \(U \subset V\) is a subspace, we can define a natural map by
\[q_U : V \to U^*, \quad x \mapsto q_U(x) = l_x\]
where \(l_x\) is just left translation by \(x\), i.e \(l_x(y) = x \cdot y\) for any \(y \in U\), then such a map is linear, and its kernel is remarkable 
\[\ker q_U = \{x \in V \mid x \cdot y = 0 , \forall y \in U\} = U^\perp\]

\begin{colorlemma}
If \((V,q)\) is a quadratic module over a field \(k\), then the following conditions are equivalent:
\begin{enumerate}[(a)]
    \item \(\mathrm{rad}(V) = 0\).
    \item rank of \(q\) is equal to \(\dim V\). 
    \item \(\mathrm{disc}(q) \neq 0\).
    \item the natural map \(q_V: V \to V^*\) is a isomorphism.
\end{enumerate}
if one of the above conditions holds, then we say that \(q\) is \textbf{non-degenerate}.
\end{colorlemma}

\begin{proof}
    \((a) \iff (b)\) follows the definition. \((a) \iff (d)\) is immediate since \(\ker q_V = \mathrm{rad}(V) = \{0\}\), so \(q_V\) is injective onto \(V^*\) which is same dimension as \(V\), hence it is an isomorphism. To finish the proof, we suppose that \(\{e_1,\dots,e_n\}\) is a basis of \(V\) and \(\{e_1^*,\dots,e_n^*\}\) is the corresponding dual basis, then \(q_V(e_i) = e_i \cdot (-)\) is just a linear functional on \(V\), hence there exists coefficients \(a_{1,i},\dots,a_{n,i}\) such that 
    \[q_V(e_i) = \sum_{k=1}^n a_{k,i} e_k^*\]
    and we can calculate the coefficients by
    \[q_V(e_i)(e_j) = e_i \cdot e_j\]
    which allows us to shows that the matrix of \(q_V\) with respect to the basis \(\{e_1,\dots,e_n\}\) and \(\{e_1^*,\dots,e_n^*\}\) is just the matrix of quadratic form \(A = (e_i \cdot e_j)_{i,j}\), hence \(\det q_V = \det A = \mathrm{disc}(q)\), which shows that \((c) \iff (d)\).
\end{proof}

\begin{example}[degenerate quadratic form]
    Given a quadratic form \(q(x,y) = x^2\) on \(k^2\), and its product is 
    \[(x_1,x_2) \cdot (y_1,y_2) = x_1y_1\]
    it is enough to check that \(\rad(V) = \mathrm{Span} \{(0,1)\}\), and so the rank of \(q\) is \(1\), and the matrix of \(q\) with respect to the standard basis is
    \[A = \begin{pmatrix}1 & 0 \\ 0 & 0\end{pmatrix}\]
    hence \(\mathrm{disc}(q) = 0\).
\end{example}

now we consider the orthogonal decomposition of quadratic forms.

\begin{definition}
    Let \((V_1,q_1)\) and \((V_2,q_2)\) be two quadratic modules over a field \(k\). The (orthogonal) direct sum is defined as \((V_1 \oplus V_2, q_1 \oplus q_2) \in \mathbf{Quad}_k\) with
    \[q_1 \oplus q_2 (u+v) = q_1(u) +q_2(v)\]
    for any \(u \in V_1\) and \(v \in V_2\). Sometimes we write \(V_1 \hat{\oplus} V_2\) to stress the direct sum of quadratic forms instead of vector spaces.
\end{definition}

it is not hard to check that the scalar product of such a quadratic form is
\[(u+v)\cdot (u'+v') = (u \cdot_1 u') + (v \cdot_2 v')\]
where \(\cdot_1\) and \(\cdot_2\) are the scalar products of \(q_1\) and \(q_2\) respectively, hence \(V_1\) and \(V_2\) are orthogonal to each other.

\begin{remark}
    If \(U\) is the supplement of \(\rad (V)\) in the sense of vector space, then \(V = U \hat{\oplus} \rad(V)\) automatically, this is owe to the definition of radical.
\end{remark}

\begin{proposition}
    Suppose that \((V,q)\) is a non-degenerate quadratic module over a field \(k\), then: 
    \begin{enumerate}
        \item Any (metric) morphism into \((V',Q')\) are injective.
        \item For any subspace \(U \subset V\), we have
        \begin{align*}
            (U^\perp)^\perp &= U \\
            \dim U + \dim U^\perp &= \dim V \\
            \rad (U) &= \rad (U^\perp) = U \cap U^\perp
        \end{align*}
        \item \(U\) is non-degenerate if and only if \(U^\perp\) is non-degenerate, and in this case, we have \(V = U \hat{\oplus} U^\perp\), or equivalently, \(U\) has a orthogonal supplement \(U^\perp\) in \(V\).
        \item If \(V = V_1 \hat{\oplus} V_2\), then \(V_1\) and \(V_2\) are non-degenerate.
    \end{enumerate}
\end{proposition}

\begin{proof}
    For (1), we assume that \(f: V \to V'\) is a morphism and \(f(x) = 0\), then by isometry
    \[0= f(x) \cdot f(y) = x \cdot y = 0\]
    for any \(y \in V\), and the condition of non-degenerate implies that \(x = 0\), hence \(f\) is injective.

    For (2), we consder the natural map \(q_U\), it induces a exact sequence
    \[0 \mapsto U^\perp \to V \xrightarrow{q_U}  U^* \to 0 \]
    so it shows that
    \[\dim  V = \dim U^\perp + \dim U^* = \dim U^\perp + \dim U \]
    replace \(U\) by \(U^\perp\) we can get
    \[\dim V = \dim U^\perp + \dim (U^\perp)^\perp\]
    a comparsion with the previous equality shows that \(\dim (U^\perp)^\perp = \dim U\), and then \(U  = (U^\perp)^\perp\) since \(U \subset (U^\perp)^\perp\) by definition. Finally, we view \(U\) as a vector space itself, and then
    \[\rad(U) = \{x \in U \subset V \mid x \cdot y = 0, \forall y \in U \subset V\} = U \cap U^\perp\]
    and \((U^\perp)^\perp = U\) shows that \(\rad(U) = \rad(U^\perp)\), which finishes the proof of (2).

    For (3), \(\rad(U) = \rad(U^\perp)\) shows the equivalence; If \(U\) is non-degenerate, then \( U \cap U^\perp = \rad U = 0\), and the sum of dimension is just \(\dim V\), which finishes the proof.

    For (4), we condsider \(q|_{V_1}\) and \(q|_{V_2}\), which gives two quadratic forms on \(V_1\) and \(V_2\) respectively, then the restriction of \(q_V\) on \(V_1\) and \(V_2\) are injective by (1), and it is enough to show that the restriction is indeed an isomorphism, so equivalently, \(V_1\) and \(V_2\) are non-degenerate.
\end{proof}

\begin{remark}
    A quadratic form \(q\) is positive if \(q(x) \geq 0\) for any \(x \in V\), we have view that \((x,y) \mapsto x \cdot y\) is a bilinear map, so if \(q\) is poistive and \(q(x) = 0\) only if \(x =0 \), then \(x \cdot y\) indeed defines a inner product on \(V\).
\end{remark}

\begin{definition}
    An element \(x\) of a quadratic module \((V,q)\) is called \textbf{isotropic} if \(q(x) = 0\); A subspace \(U \subset V\) is called \textbf{isotropic} if all elements are isotropic.
\end{definition}

Notice that in inner product space, the isotropic element is just the zero vector, so in general, such a vector is the special role for quadratic forms theory.

\begin{remark}
    By definition, for a subspace \(U \subset V\) we have
    \[2 x \cdot y = q(x+y) -q(x) -q(y)\]
    for any \(x, y \in V\). If \(U\) is isotropic, and we set \(x,y \in U\), then \(x +y \in U\), and then \(x \cdot y = 0\), which is enough to shows that \(U \subset U^\perp\); Conversely we set \(x = y\), then
    \[q(x) = x \cdot x = 0\]
    for any \(x \in U\), which shows that \(U\) is isotropic. Hence we have the equivalence
    \[U \text{ is isotropic } \iff U \subset U^\perp\]
\end{remark}


\begin{definition}
    A hyperbolic plane is defined as a two-dimensional quadratic module \((V,q)\) with an isotropic-basis \(\{x,y\}\) such that \(x \cdot y \neq 0\).
\end{definition}

\begin{example}
    Let \(q(x,y) = xy\) be a quadratic form on \(k^2\), then \(\{(1,0),(0,1)\}\) is an isotropic-basis since \(q(1,0) = q(0,1) = 0\), and \[(1,0) \cdot (0,1) = \frac{1}{2}(q(1,1)-q(1,0)-q(0,1)) = \frac{1}{2}(1-0-0) = \frac{1}{2} \neq 0\]
    so the matrix of quadratic form is 
    \[\begin{pmatrix}
        0 & 1/2 \\
        1/2 & 1 
    \end{pmatrix}\]
\end{example}

More generally, each matrix of the hyperbolic plane is of the form 
\[\begin{pmatrix}
    0 & 1 \\
    1 & 0
\end{pmatrix}\]
after a nomalization of the basis (change \(x,y \) to be \(x, \frac{y}{x \cdot y}\)).

\begin{colorproposition} \label{translation}
    Let \((V,q)\) be a quadratic module over a field \(k\) containing a non-zero isotropic element \(x\), then
    \begin{enumerate}
        \item There exists a hyperbolica plane \(H\) of \(V\) such taht \(x \in H\).
        \item \(q\) is surjective onto \(k\).
    \end{enumerate}
\end{colorproposition}

\begin{proof}
    The condition of non-degenerate allows us to find a \(z \in V\) such that \(x \cdot z  \neq 0 \), and we set
    \[y = 2z + ax ,  \quad a = - \frac{z \cdot z}{x \cdot z}\]
    then we can verify that \(y\) is isotropic and \(x \cdot y \neq 0\), which shows that \(H = \mathrm{Span}_k \{x ,y\}\) is a hyperbolic plane containing \(x\). 
    
    (2) is the corollary of (1), it is enough to show that \(q|_H\) is surjective onto \(k\). We suppose that \(\{x,y\}\) is a isotropic-basis for H, then for any \(a \in k\), we can find \(b = \frac{a}{2 x \cdot y} \in k\) such that \(q(x +by)= a\), which finishes the proof.
\end{proof}

\textcolor{blue}{Finally, we generalize the orthonormal basis of inner product space to the context of quadratic forms, the existence of such a basis can be guarantedd by the Gram-Schmidt process in the context of inner product space, and same work can done in spite of some technical details.}

\begin{definition}
    A basis \(\{e_1,\dots,e_n\}\) of a quadratic module \((V,q) \in \mathbf{Quad}_R\) is called \textbf{orthogonal} if \(e_i \cdot e_j = 0\) for any \(i \neq j\) such that
    \[V = Re_1 \hat{\oplus} Re_2 \hat{\oplus} \cdots \hat{\oplus} Re_n\]
    or equivalently, the matrix of \(q\) with respect to the basis is diagonal \(A = \mathrm{diag}(a_1,\dots,a_n)\) such that 
    \[q(x_1e_1+ \dots + x_n e_n) = a_1x_1^2+ \dots+ a_nx_n^2\]
\end{definition}

\begin{colortheorem}
    Any quadratic module \((V,q)\) over a field \(k\), then
    \begin{enumerate}
        \item \(V\) admits an orthogonal basis.
        \item If \(V\) is non-degenerate and of dimension \(n \geq 3\), then for any two orthogonal bases \(\mathcal{E} = \{e_1,\dots,e_n\} \) and \(\mathcal{F} = \{f_1,\dots,f_n\}\), there exists a finite sequence of orthogonal bases \(\mathcal{E}_0,\dots,\mathcal{E}_m\) such that \(\mathcal{E}_0 = \mathcal{E}\), \(\mathcal{E}_m = \mathcal{F}\), and \(\mathcal{E}_i\) is contiguous with \(\mathcal{E}_{i+1}\) for any \(0 \leq i \leq m-1\).
    \end{enumerate}  
\end{colortheorem}

\subsection{Witt's theorem}

As we have defined, a metric morphism between two quadratic modules is a linear map preserving the quadratic form, a useful result is that such a morphism on the subspace can be extended to the whole space.

\begin{colortheorem}[Witt's theorem]
    Let \((V,q)\) and \((V',q')\) be two \textbf{non-degenerate} quadratic modules over \(k\), and \(U\) is a subspace of \(V\) with a morphism
    \[s: U \to V'\]
    \begin{enumerate}
        \item If \(U\) is degenerate, then \(s\) can be extended to a morphism on \(U_1\) such that \(U \subset U_1\).
        \item If \(U\) is non-degenerate and two quadratic modules are isomorphic, then \(s\) can be extended to a morphism on whole \(V\).
    \end{enumerate}
\end{colortheorem}

\begin{proof}
    \begin{itemize}[label=$\blacktriangleright$]
        \item   For the first case, we take \(x \in \rad U - \{0\}\) and \(l \in U^*\) such that \(l(x) = 1\). \(V\) is non-degenerate, so \(q_V: V \to V^*\) is an isomorphism, and we extend \(l\) to a linear functional on \(V\), then there exists a unique \(y \in V\) such that \(l(u) = u \cdot y\),and we set \(U_1 = ky \hat{\oplus} U\). For \(U' = s(U)\) with \(x' =  s(x)\), we consider \(l' = l \circ s^{-1}\) where \(s^{-1}\) is the inverse of injection \(s\) on \(U\), similarly there exists a unique \(y' \in V'\) such that \(l'(u') = u' \cdot y'\), and then we can construct a linear map \(s_1: U_1 \to U_1' = U' \hat{\oplus} k y'\) by
    \[s_1(u+ cy) = s(u) + cy'\]
    and we can verify that \(s_1\) is indeed a metric morphism with an assumption that \(y \cdot y = y' \cdot y' = 0\)
    \begin{align*}
        (s(u) + cy') \cdot (s(u) + cy') &= s(u) \cdot s(u) + 2c \, l'(s(u))+ 0 \\
        &= u \cdot u + 2 c \, l \circ s^{-1} \circ s (u) + c^2 y \cdot y \\
        &= (u + cy) \cdot (u + cy)
    \end{align*}    

    The assumption can be always satisfied as following: If \(y \cdot y \neq 0\), we consdier \(y_0 = y - \frac{y \cdot y}{2} x \) and then by \(x \in \rad U\)
    \[u\cdot y_0= u \cdot y - 0  = l(x)\]
    and 
    \[y_0 \cdot y_0 = y \cdot y -(y\cdot y )l(x) + 0  = 0\]
    so we can replace \(y \) by \(y_0\), and same work can be done for \(
    y'\).

    \item For The second case, we prove it by induction on the dimension of \(U\). If \(\dim U = 1\) and \(x \in U\) is a non-zero vector, then we set \(y = s(x)\) and \(y \cdot y = x \cdot x\) by isometry. We can claim that \(x + y \) and \(x - y\) can not be all isotropic, otherwise by expansion we have
    \[2 x \cdot x \pm 2x \cdot y = 0\]
    which shows that \(x \cdot x  = 0\) which contradicts the non-degenerate condition (\(\rad U = \{0\}\)). Hence we fix \(e  = \pm 1\) such that \(z = x + ey \) is not isotropic, and then subspace \(\mathrm{Span}_k\{z\}\) is non-degenerate, so it admits an orthogonal supplement \(H\) such that \(V = kz \hat{\oplus} H\), and we can construct a automorphism \(\sigma: V \to V'\) by
    \[\forall a = cz + h \in V , \sigma(cz+h) = -cz + h\]
    which implies \(\sigma(x ) = -ey\), so \(-e \cdot \sigma\) can be seen as an extension of \(s\).

    If \(\dim U \geq 2\), there exists a decomposition \(U = U_1 \hat{\oplus} U_2\), and then we consider \(s|_{U_1}\), by induction there exists an extension \(\sigma_1: V  \to V'\) such that \(\sigma_1|_{U_1} = s|_{U_1}\). By composition \(\sigma_1^{-1}\circ s : U \to V\), we can find such a map is identity on \(U_1\), and its restriction on \(U_2\) can be extended to a morphism \(\tilde{\sigma}: V \to V\) by induction such that \(\tilde{\sigma}\) is still identity on \(U_2\), and then we consider \(\sigma_1 \circ \tilde{\sigma} : V \to V'\): for any \(s_1 \in U_1\), we have 
    \[\sigma_1 \circ \tilde{\sigma} (s_1) = \sigma_1(s_1) = s(s_1)\]
    and for any \(s_2 \in U_2\) we have
    \[\sigma_1 \circ \tilde{\sigma} (s_2) = \sigma_1 \circ \sigma_1^{-1} \circ s(s_2) = s(s_2) \]
    which gives a desired extension of \(s\) on \(V\).
      \end{itemize}
\end{proof}

\subsection{Translation}

To do the classification of quadraic form, we need to write it explicitly, i.e. 
\[f(X) = \sum_{i} a_{i,i}X_i^2+ 2\sum_{i < j} a_{i,j}X_iX_j\]
denotes a quadratic form over \(k\) with \(n\) variables, and we introduce some notation to discuss.

\begin{definition}
    The matrix of the quadratic form \(f\) is defined as 
    \[A = (a_{i,j})_{i,j}\]
    two quadratic forms \(f\) and \(g\) are equivalent if the two matrices of them are similiar, and we denote it by \(f \sim g\).
\end{definition}

\begin{remark}
    We consider the hyperbolic plane \(f(X_1,X_2) = X_1X_2\), another similar form is \(g(X_1,X_2) = X_1^2 - X_2^2\), to see that they are equivalent, we set \(Y_1 = X_1+X_2\) and \(Y_2 = X_1 - X_2\), then we have
    \[g(X_1,X_2) = Y_1Y_2\]
\end{remark}

\begin{definition}
    For a quadratic form \(f\) over \(k\) with two variables, it is called hyperbolic if \(f \sim X_1X_2\), or equivalently, quadratic modules \((k^2,f)\) is a hyperbolic plane.
\end{definition}

\begin{definition}
    Let \(f(X_1,\dots,X_n) \) and \(g(X_1,...,X_m)\) be two quadratic forms over \(k\), then \(f \pm g\) is defined as a quadratic form over \(k\) with \(n+m\) variables by
    \[(f \pm g)(Y_1,\dots,Y_{n+m}) = f(Y_1,\dots,Y_n) \pm g(Y_{n+1},\dots,Y_{n+m})\]
\end{definition}

such a definition is equivalent to the orthogonal direct sum of quadratic forms.

\begin{definition}
    A quadratic form \(f\) over \(k\) represents an element \(a \in k\) if there exists a vector \(x \in k^n\) such that \(f(x) = a\).
\end{definition}

so for a quadratic equation, \(x\) is a solution if and only if \(f\) represents \(0\), and actually, isotropic elements are just the solutions of the quadratic equation \(f(x) = 0\).

\begin{proposition}
    If \(f\) is a non-degenerate quadratic form over \(k\) with a non-trival zero, then \(f \sim h + f'\) with \(h\) a hyperbolic plane, and \(f\) represents all elements of \(k\).
\end{proposition}

It is just a translation of Proposition~\ref{translation}, nothing new should be added.

\begin{colorcorollary}
    Let \(g(X_1,\dots,X_{n-1})\) be a non-degenerate quadratic form over \(k\), and \(a \in k^*\) non-zero, then the following conditions are equivalent:
    \begin{enumerate}
        \item \(g\) represents \(a\).
        \item \(g \sim h(X_1,\dots,X_{n-2}) + a Z^2\).
        \item \(f = g - aZ^2\) represents 0.
    \end{enumerate}
\end{colorcorollary}

\begin{proof}
    (1) implies (2): \(x \cdot x = g(x) = a \neq 0\), so \(kx\) is a non-degenerate subspace of \(k^{n-1}\), which admits an orthogonal supplement \(H\), hence \(V = H \hat{\oplus} kx\) allows us to write such a equivalent.

    (2) implies (3): we can write 
    \[g(Y_1,\dots,Y_{n-1}) = h(X_1,\dots,X_{n-2}) + a Z^2\]
    where \((X_1,\dots,X_{n-2},Z)\) is a change basis of \((Y_1,\dots,Y_{n-1})\), then we can write
    \[f = g-aZ^2 = h(X_1,\dots,X_{n-2}) + aX_{n-1}^2 - aX_n^2\]
    and a non-trival zero \((0,\dots,0,1,1)\) can be determined.

    (3) implies (1): If \((x_1,\dots,x_{n-1},z)\) is a non-trival solution, we consider two cases: when \(z = 0\), then \(g(x_1,\dots,x_{n-1}) = 0\) and \(g\) represents \(0\); when \(z \neq 0\), then \((x_1/z,\dots,x_{n-1}/z)\) is a solution of \(g - a =0\).
\end{proof}

\begin{colortheorem}
    Each quadratic form \(f\) over \(k\) with \(n\) variables is equivalent to a diagonal form \(f \sim a_1X_1^2 + \cdots + a_n X_n^2\).
\end{colortheorem}

\begin{proof}
    
\end{proof}




\newpage


\subsection{Quadratic forms over \texorpdfstring{the $p$-adic field}{the p-adic field}}

We firstly consider the case over finite field, then it can be lifted to the \(p\)-adic field, and we always assume that the characteristic of \(k\) is not two.

\begin{proposition}
    A quadratic form \(q\) over \(\ff_q\) of rank \(\geq 2\) (resp. of rank \(\geq 3\)) represents all elements of \(\ff_q^*\) (resp. \(\ff_q\)).
\end{proposition}

\newpage

\section{Modular Form}

\subsection{Modular Group }

\begin{definition}
    \textbf{Modular group} is 
    \[\mathrm{SL}_2(\zz) = \left\{ \begin{pmatrix} a & b \\ c & d \end{pmatrix} \in \mathrm{GL}_2(\zz) : ad - bc = 1 \right\}\]
    and for any \(N \in \nn\), \textbf{the principal congruence subgroup of level \(N\) }is defined as
    \[\Gamma(N) = \{M \in \mathrm{SL}_2(\zz) : M \equiv I_2 \mod N\}\]
    A group \(\Gamma\) is a  \textbf{congurence subgroup} if there exists \(N \in \nn\) such that \(\Gamma(N) < \Gamma < \mathrm{SL}_2(\zz)\).
\end{definition}

It is not hard to check that \(\Gamma(N)\) is a subgroup of modular group, and in particular \(\Gamma(1)\) is just the modular group.

\begin{example}
    For any \(N \in \nn\), we can define
    \[\Gamma_0(N) = \{\begin{pmatrix}
    a & b \\ c & d 
    \end{pmatrix} \in \mathrm{SL}_2(\zz) : c \equiv 0 \mod N\}\]
    this is an important congurence sungroup in Wile's proof of FLT, and another congurence sungroup is
    \[\Gamma_1(N) = \{\begin{pmatrix}
    a & b \\ c & d 
    \end{pmatrix} \in \mathrm{SL}_2(\zz) : a \equiv d \equiv 1 \mod N, c \equiv 0 \mod N\}\]
    and it is not hard to check that \(\Gamma(N) < \Gamma_1(N) < \Gamma_0(N) < \mathrm{SL}_2(\zz)\).
\end{example}


\begin{definition}
    The left action of \(\mathrm{SL}_2(\zz)\) on the upper half plane \(\hh\) is defined by Mobius transformation
\[g\cdot z =  \frac{az+b}{cz+d}\]
for any \(g = \begin{pmatrix} a & b \\ c & d \end{pmatrix} \in \mathrm{SL}_2(\zz)\) and \(z \in \hh\), and we define \textbf{modular cocycle }by
\[j(g,z) = cz+d\]
\end{definition}

For any \(M \in \mathrm{SL}_2(\zz)\), we can check that
\[\mathrm{Im}(M\cdot z ) = \frac{\det M \cdot \mathrm{Im}(z)}{|cz+d|^2}\]
and it is not difficult to check such an action is well-defined.

\begin{lemma}
    For any \(g,h \in \mathrm{SL}_2(\zz)\) and \(z \in \hh\),
    
\end{lemma}

\section{Local field}

In this section, all ring will be assumed to be commutative with identity, and it has no zero divisor, i.e. each ring is an integral domain, so this section is a specific case of commutative algebra.
\subsection{Local ring and Valuation ring}

\begin{definition}
    A local ring is an integral domain \(R\) with a unique maximal ideal \(\mathfrak{m}\), and we denote it by \((R,\mathfrak{m})\), and the field \(k = R/\mathfrak{m}\) is called the residue field of \(R\).
\end{definition}

A fact about local rings is that 
\[R^\times = R - \mathfrak{m}\]
all elements in the maximal ideal are non-invertible, which gives an equivalent definition of local ring: \textcolor{blue}{all non-units elements form an ideal}.

\begin{lemma}[Nakayama]\label{Nakayama}
    Let \((R,\mathfrak{m})\) be a local ring, and \(M\) be a finitely generated \(R\)-module, if \(\mathfrak{m}M = M\), then \(M = 0\).
\end{lemma}

It is a special case of theorem \ref{Nakayama}, but a short proof can be done based on the properties of local ring.

\begin{proof}

\end{proof}


\begin{corollary}
    \(\mathfrak{m}/\mathfrak{m}^2\) is a \(k\)-vector space, and it is finite-dimensional if \(R\) is \textbf{noetherian}, then the dimension is equal to the minimal number of generators of \(R\)-module  \(\mathfrak{m}\).
\end{corollary}

\begin{proof}
    We define the scalar multiplication of \(k\) on \(\mathfrak{m}/\mathfrak{m}^2\) by
    \[(a+ \mathfrak{m}) \cdot (x + \mathfrak{m}^2) = ax + \mathfrak{m}^2\]
    with \(a \in R\) and \(x \in \mathfrak{m}\), it is not hard to check that it is well-defined. If \(R\) is noetherian, then \(\mm\) is finitely generated by finite elements \(\{x_1,\dots,x_n\}\), so any element \(x \in \mm\) can be written as
    \[x = \sum_i c_i \cdot x_i + \sum_j b_j x_j\]
    with \(c_i \in k\) and \(b_j \in \mm\), hence mod \(\mm^2\) we have
    \[x + \mm^2 = \sum_i c_i \cdot (x_i + \mm^2)\]
    which shows that \(\dim_k(\mm/\mm^2) \leq n < +\infty\). Now we suppose that \(n\) is the smallest number of generators of \(\mm\), and \(N\) is the submodule of \(M\) generated by \(\{x_1,\dots,x_n\}\), then \(\mm = N + \mm^2\), hence by Nakayama lemma, we have \(\mm = N\), which shows that \(\dim_k(\mm/\mm^2) = n\), then 
    \[\mm = N + \mm^2\]
    we can apply Nakayama lemma to \(R\)-module \(M = \mm/N\) such that
    \[\mathfrak{m}(\mm/N) = (\mm^2+N)/N = \mm/N\]
    which shows that \(M = \mm/N = 0\), hence \(\mm = N\), which shows that \(\dim_k(\mm/\mm^2) = n\).
\end{proof}

\begin{remark}
    The result will make sense in algebraic geometry to define the dimension of a variety at a local point.
\end{remark}


\begin{definition}
    A valuation ring is a ring \(R\) such that for any \(x \in K = \mathrm{Frac}R\), either \(x \in R\) or \(x^{-1} \in R\).
\end{definition}

\begin{remark}
    A valuation ring is a local ring, to show that, it is enough to show that the set of non-units forms an ideal.
\end{remark}

\begin{remark}
    The valuation ring is integral closed (over \(K\)): assume \(x \notin R\) is integral over \(K\), then there exists \(a_i \in R\) such that
    \[x^n + a_{n-1}x^{n-1} + \cdots + a_1x + a_0  = 0\]
    by definition \(x ^{-1} \in R\), so we can conclude that
    \[x = a_{n-1} + a_{n-2}x^{-1} + \cdots + a_0 x^{1-n} \in R\]
    which leads to a contradiction, so \(\bar{R} = R\) in \(K\).
\end{remark}




\begin{proof}
    
\end{proof}

\begin{definition}
    Let \((\Gamma,+,<)\) be a totally ordered group, and \(K\) be a field, then a valuation of \(K\) is a map \(v: K \to \Gamma \cup \{\infty\}\) such that\\
    (1) \(v(x) = \infty \iff x = 0\).\\
    (2) \(v(xy) = v(x) + v(y)\).\\
    (3) \(v(x+y) \geq \min\{v(x),v(y)\}\).\\
    If \(\Gamma = \rr\), the valuation is called \textbf{height one valuation}; If \(v(K^*) = \zz\), then it is called a \textbf{discrete valuation}.
\end{definition}

\begin{remark}
    If \(v\) is a valuation on field \(K\), and we define a set
    \[R_v : = \{x \in K : v(x) \geq 0\}\]
    with \(0\) the identity of \(\Gamma\), then \(R\) is a valuation ring with the unique maximal ideal
    \[\mm_v  = \{x \in K: v(x) > 0\}\]
    and the unit group is \[R_v^\times = \{x \in K: v(x) = 0\}\].
    Conversely, for any valuation ring \(R\) with fraction field \(K\), we define \((\Gamma,+) = (K^\times/R^\times, \times)\), and it is a totally ordered group with the order defined by
    \[\bar{x} \leq \bar{y} \iff xy^{-1} \in R\]
    then the quotient map \(\pi: K^\times \to \Gamma\) can be extended to a valuation by defining \(\pi(0) = \infty\). The proof can be left as an exercise, and it gives a surjection \(v \mapsto R_v\) between the set of valuations on \(K\) and the set of valuation rings of \(K\).

\end{remark}

We will restrict our attention to the discrete valuation ring (DVR), i.e. the valuation ring with a discrete valuation, which is a special case and has many good properties to consider.

\begin{proposition}
    A integral domain \(R\) is a discrete valuation ring if and only if \(R\) is a principal ideal domain with a unique non-zero maximal ideal \(\mm\).
\end{proposition}
\begin{proof}
    If \(R\) is a PID, then any maxiaml ideal is generated by a irreducible element, if such a element is unique and denoted by \(\pi\), then any non-zero element is of the form \(\pi^n u\) where \(u\) is a unit, which allows us to define a valuation \(v = n\) with respect to the exponent of \(\pi\), and such a valuation is clearly discrete.

    Conversely, valuation ring is a local ring with a unique maximal ideal. There exists a element \(\pi \in R\) such that \(v(\pi) = 1>0\), so \(\pi \in \mm\) and it is not unit. Then for any \(x \in \mm\), \(x/\pi \in R\) since \(v(x/\pi) \geq 0\), which shows that \(x \in (\pi)\) and \(\mm = (\pi)\), which shows that such a maximal ideal is no-zero. Finally, we prove that \(R\) is a PID. For any non-zero ideal \(I\) we can suppose that \(m\) is the minimum of \(v(I)\) in \(\nn\), so there exists \(a \in I\) of the form \(a = \pi^m v\) for some unit \(v\), which is enough to finish the proof.
\end{proof}

\begin{remark}
    Equivalently, the condition can be replaced by: \(R\) is a local and principal, and \(R\) is not a field. In serre's book, another equivalent condition is that \(R\) is a integral closed noetherian integral domain with a unique non-zero maximal ideal \(\mm\).
\end{remark}

\subsection{Dedekind domain}

\begin{definition}
    A Dedekind domain is a ring \(R\) such that (1) noetherian integral domain; (2) integral closed; (3) dimension one, i.e. any non-zero prime ideal is maximal.
\end{definition}

\begin{remark}
    Localization keeps three properties of Dedekind domain, so the localization of a Dedekind domain at a prime ideal is still a Dedekind domain, here we discuss (3).

    If \(\mathfrak{p}\) is a non-zero prime ideal of \(R\), then there exists a correspondence:
    \[\{\mathrm{Spec}(R_\mathfrak{p})\} \leftrightarrows \{ \mathfrak{q} \in \mathrm{Spec}(R): \mathfrak{q}  \subset \mathfrak{p}\}\]
    the one-dimensional condition shows that \(\mathfrak{q}\) is zero or \(\mathfrak{p}\), so clearly the local ring \(R_\mathfrak{p}\) is one-dimensional.
\end{remark}

Review the definition of Dedekind domain here, a typical example is the ring of integers of a number field, in number theory we often consider the localization of such ring at a prime ideal, which is usually a DVR or more precisely, a extension of \(\qq_p\).

\begin{proposition}
    A ring \(R\) is a DVR if and only if it is a local Dedekind domain (a Dedekind domain with a unique non-zero prime ideal, so it is a local ring).
\end{proposition}

\begin{proof}
    It is enough to prove the necessity.
\end{proof}

\begin{corollary}
    If \(R\) is a noethrian integral domain, then \(R\) is a Dedekind domain if and only if \(R_\mathfrak{p}\) is a DVR for any non-zero prime ideal \(\mathfrak{p}\).
\end{corollary}







\subsection{Absolute value and completion}

To talk about the topology of a valuation ring and its fraction field, we tend to define a metric on it, to do that, notice that exponential map is a natural bijection between \(\rr\) and \(\rr_{>0}\).

\begin{definition}
    A normed (or valued) field is a field \(K\) together with a positive function \(|\cdot|: K \to \rr_{\geq 0}\) such that\\
    (1) \(|x| = 0 \iff x = 0\).\\
    (2) \(|xy| = |x||y|\).\\
    (3) \(|x+y| \leq |x| + |y|\).\\
    furthermore, if (3) is replaced by a stronger condition \(|x+y| \leq \max\{|x|,|y|\}\), then it is called a non-archimedean absolute value.
\end{definition}

for the height one valuation, we can easily verify the following correspondence 
\[\{\text{height one valuations on } K \} \longleftrightarrow  \{\text{non-archimedean absolute values on } K\}\]
by
\[v \mapsto |\cdot|_v = e^{-v(\cdot)}\]

The absolute value gives a metric to the field by
\[d(x,y) = |x-y|\]
so as a metric space, it is natural to consider the completion of the field as real number, so a natural definition can be given here.

\begin{definition}
    Let \((K,|\cdot|_K)\) be a normed field, the completion of \(K\) is a complete normed field \((L,|\cdot|_L)\) together with an isometric embedding \(\iota: K \to L\) such that (i) \(|\iota(x)|_L = |x|_K\) for any \(x \in K\); (ii) \(\iota(K)\) is dense in \(L\).
\end{definition}

By definition, an example is that \(\rr\) is a completion of \(\qq\), while \(\cc\) is not a completion of \(\qq\) although it is complete. In the view of grammar, we should say "the" completion instead of "a" completion, which can be viewed as a consequence of the following universal property of completion:

\begin{proposition}[UP-completion] \label{UP-completion}
    Let \((K,|\cdot|_K)\) be a normed field with a completion \((L,|\cdot|_L)\) and an isometric embedding \(\iota: K \to L\), then for any \textbf{complete} normed field \((E,|\cdot|_E)\) with a isometric field homomorphism \(f:K \to E\), the morphism \(f\) can be uniquely extended to an (injective) isometric morphism \(\tilde{f}:L \to E\) such that the following diagram commutes:
    % https://q.uiver.app/#q=WzAsMyxbMCwwLCJLIl0sWzEsMCwiTCJdLFswLDEsIkUiXSxbMCwxLCJcXGlvdGEiLDAseyJzdHlsZSI6eyJ0YWlsIjp7Im5hbWUiOiJob29rIiwic2lkZSI6InRvcCJ9fX1dLFswLDIsImYiLDJdLFsxLDIsIlxcdGlsZGV7Zn0iLDAseyJzdHlsZSI6eyJib2R5Ijp7Im5hbWUiOiJkYXNoZWQifX19XV0=
\[\begin{tikzcd}
	K & L \\
	E
	\arrow["\iota", hook, from=1-1, to=1-2]
	\arrow["f"', from=1-1, to=2-1]
	\arrow["{\tilde{f}}", dashed, from=1-2, to=2-1]
\end{tikzcd}\]
\end{proposition}

\begin{proof}
    
\end{proof}

\begin{corollary}
    Each normed field has a unique completion up to isometric isomorphism.
\end{corollary}

\begin{proof}
    
\end{proof}

We first consider the archimedean absolute value, this is a classic and rigid example, and we conclude here.

\begin{lemma}
    Let \((K,|\cdot|_K)\) be a normed field, then the following conditions are equivalent:\\
    (1) It does not satisfy the stronger triangle inequality, i.e. there exists \(x,y \in K\) such that \(|x+y|_K > \max\{|x|_K,|y|_K\}\).\\
    (2) There exists a integer \(N \in \nn\) such that \(|N\cdot 1_K|>1\).\\
    (3) The set \(\{|n \cdot 1_K| : n \in \nn\}\) is unbounded.
    If one of the above conditions holds, then \((K,|\cdot|_K)\) is called an \textbf{archimedean} normed field.
\end{lemma}
\begin{proof}
    
\end{proof}

\begin{proposition}
    The complete archimedean normed field is isometric to \(\rr\) or \(\cc\) with the usual absolute value.
\end{proposition}

\begin{proof}
    Firstly, archimedean field must be of characteristic zero, otherwise 
    \[\{|n \cdot 1|_K : n \in \zz\} = \{0,1,\dots,p-1\}\]
    with the assumpion that \(\mathrm{char} K = p \neq 0\). Then we can embed \((\qq,|\cdot|_{\infty})\) into \(K\) by \(1 \mapsto 1_K\), and then the universal property of completion (proposition \ref{UP-completion}) shows that there exists an injective morphism from \(\rr\) to \(K\), which shows that \(K\) must be the real commutative division algebra, hence \(K \cong \rr\) or \(K \cong \cc\).
\end{proof}

We have seen that the height one valuation is just the non-archimedean absolute value, so for any non-archimedean normed field \((K,|\cdot|_K)\), we usually use the notation
\[\mathcal{O}_K = \{x \in K: |x|_K \leq 1\} = \{x \in K : v(x) \geq 0\}\]
to denote the valuation ring of \(K\), and it is a local ring with the unique maximal ideal
\[\mm = \{x \in K : |x|_K <1 \} = \{x \in K : v(x) > 0 \}\]
with the residue field
\[k = \mathcal{O}_K/\mm\]
such the convention is widely used in local field. As a special metric space, non-archimedean normed field can be generalized to be a ultra-metric space, which is a metric space \((X,d)\) such that the metric \(d\) satisfies a stronger condition
\[d(x,y) \leq \max\{d(x,z),d(z,y)\}\]
for any \(x,y,z \in X\). Hence, we can study its analytical properties, and we can find the sequence in such space will satisfy many stronger properties.

\begin{lemma}
    Let \((X,d)\) be a ultra-metric space, then 
    \begin{enumerate}
        \item For any \(y \in B(x,r)\), we have \(B(x,r) = B(y,r)\).
        \item For any \(y \in \overline{B(x,r)}\), we have \(\overline{B(x,r)} = \overline{B(y,r)}\).
        \item \(B(x,r)\) and \(\overline{B(x,r)}\) are both closed and open, consequently, \(X\) is totally disconnected.
        \item A sequence \(\{x_n\}\) is Cauchy if and only if \(\lim_{n\to \infty} d(x_n,x_{n+1}) = 0\).
    \end{enumerate}
\end{lemma}

The proof can be found in any textbook about p-adic analysis, we omit it here. The last of section is to introduce some sepcific examples of non-archimedean normed fields.

\subsection*{Example: p-adic number field}

There are many ways to construct such a field, a classic way is to see it as a different completion of \(\qq\).

\begin{definition}
    Let \(p\) be a prime number, then the \(p\)-adic valuation on \(\zz\) is defined by
    \[v_p(n) = \max \{m \in \nn: p^m \mid n \}\]
    such a definition can be extended to \(\qq\) by 
    \[v_p\left(\frac{a}{b}\right) = v_p(a) - v_p(b)\]
    for any \(a,b \in \zz\) with \(b \neq 0\), and we convention that \(v_p(0) = \infty\). 
\end{definition}
It is not hard to verify that \(v_p\) is a height one valuation on \(\qq\), and then we can define the \(p\)-adic absolute value by
\[|x|_p = e^{-v_p(x)}\]
such that \((\qq,|\cdot|_p)\) is a non-archimedean normed field, and its valuation ring is
\[\zz_{(p)} = \{\frac{a}{b} \in \qq : p \neq b\}\]
it can be viewed as the localization of \(\zz\) at the prime ideal \((p)\). Such a normed field and its valuation ring are not complete.

\begin{example}
    We consider a sequence defined by
    \[x_n = \sum_{k = 1}^n p^{k }\]
    and then 
    \[|x_{n+1}-x_n|_p = |p^{(n+1)}|_p = e^{-n-1}\]
    shows that such a sequence is Cauchy, but its limit does not exist in \(\qq\).
\end{example}

\subsection*{Example: functional field}

Let \(k\) be a field, then we define a valuation on \(k[t]\) by
\[v_t(P) := \max \{k \in \nn: t^k \mid P(t)\}\]
for any \(P \in k[t]\), such a valuation can be extended to \(k(t)\) by
\[v_t\left(\frac{P}{Q}\right) = v_t(P) - v_t(Q)\]
for any \(P,Q \in k[t]\) with \(Q \neq 0\), and we convention that \(v_t(0) = \infty\). It is not hard to verify that \(v_t\) is a discrete valuation on \(k(t)\).

\begin{example}
    If we define the \textbf{t-adic absolute value} by
    \[|f|_t : = e^{-v_t(f)}\]
    then \(k(t)\) is not necessary complete, to see that we consider a sequence
    \[f_n(t) = \sum_{k=0}^n t^k\]
    and then 
    \[|f_{n+1}-f_n|_t = e^{-n-1}\]
    shows that such a sequence is Cauchy, while its limit is a formal power series, it makes sense in \(k((t))\), i.e. \textbf{the field of Laurent series}, which contains all elements of the form
    \[P(t) = \sum_{k \geq N}^{+\infty}a_k t^k\]
    for some \(N \in \zz\).
\end{example}


\begin{proposition}
    \((k((t)),|\cdot|_t)\) is a complete non-archimedean normed field, and its valuation ring is
    \[\mathcal{O}_t = k[t]_{(t)} = \{\frac{P}{Q} \in k(t) : Q(0) \neq 0\}\]
    and the residue field is \(k\).
\end{proposition}

\begin{proof}
    
\end{proof}

Another natural absolute value is given by degree function, we recall that degree function on a polynomial ring \(k[t]\) satisifes
\begin{enumerate}
    \item \(\deg P = -\infty \iff P = 0 \).
    \item \(\deg(PQ) = \deg P + \deg Q\).
    \item \(\deg(P+Q) \leq \max\{\deg P,\deg Q\}\)
\end{enumerate}
so it is enough to define an valution on \(k(t)\) by
\[v_\infty(P/Q) = \deg Q - \deg P\]
with a convention that \(v_\infty(0) = \infty\), and then we can define a non-archimedean absolute value by
\[|f|_{\infty} = e^{-v_{\infty}(f)}\]

\begin{proposition}
    \(k((1/t))\) is the completion of \(k(t)\) with respect to the absolute value \(|\cdot|_\infty\), and its valuation ring is
    \[\mathcal{O}_\infty = k[1/t]_{(1/t)}\]
    and its residue field is \(k\). 
\end{proposition}

\begin{proof}
    
\end{proof}


\subsection{Hensel's lemma}

\begin{example}
    \(\zz_p\) is a different ring from \(\zz\), for example \(\sqrt{2} \in \zz_7\). We assume an element of \(\zz_7\) with 7-digit expansion
    \[a = a_0 + a_1 7 +a_2 7^2 + \cdots\]
    with \(a_i \in \{0,1,\cdots,6\}\) such that \(a^2 = 2\), then modulo 7 we have
    \[2 = a^2 = a_0^2 \mod 7    \]
    it has two solution \(a_0 = 3,4\). We fix \(a_0 =3 \), and then
    \[2 = a^2 = 9 + 6a_1 \cdot 7 + \cdots\]
    similarly, modulo \(7^2\) we have
    \[2 = 9 +42 a_1 \mod 49\]
    it has a unique solution \(a_1 = 6\), and then we can continue to find each \(a_i\) to get
    \[\sqrt{2} = (\cdots,a_2,a_1,a_0)_7 = \cdots 16213_7\]
    or we lift the solution starting from \(a_0 = 4\), and we can get another solution
    \[-\sqrt{2} = \cdots 50454\]
\end{example}

\begin{theorem}[Hensel's lemma] 
    Let \((K,|\cdot|)\) be a complete non-archimedean normed field, and we assume that \(f \in \mathcal{O}_K[X]\) is a polynomial such that 
    \[|f(a_0)| < |f'(a_0)|^2\]
    for some \(a_0 \in \mathcal{O}_K\), then there exists unique \(a \in \mathcal{O}_K\) such that \(f(a ) = 0\) and\\
    (1) Newton's sequence starting from \(a_0\)
    \[a_{n+1} = a_n - \frac{f(a_n)}{f'(a_n)}\]
    quadratically converges to \(a\).\\
    (2) \(|a - a_0|= |f(a_0)/f'(a_0)| < |f'(a_0)|\).\\
    (3)\(|f'(a)| = |f'(a_0)|\).
\end{theorem}

\begin{proof}
    
\end{proof}

\begin{corollary}[simple root]
    In above setting, if 
    \[f(a_0) = 0 \mod \mathfrak{m}, \quad f'(a_0) \neq 0 \mod \mathfrak{m}\] 
    then same conclusion holds, a root \(a \in \mathcal{O}_K\) can be lifted as the root of \(f\) and \(a = a_0 \mod \mathfrak{m}\). 
\end{corollary}

\begin{proof}
    
\end{proof}

\begin{corollary}[quadratic residue]
    For any non-zero element \(u = p^kv\) with \(v \in \zz_p^\times\) and \(k \in \zz\), then \(u\) is a square if and only if 
    \begin{itemize}
        \item when \(p =2\), \(k \) is even and \(v =1 \mod 8\).    
        \item when \(p\) is odd, \(k\) is even and \(v\) is a square modulo \(p\).
    \end{itemize}
    As a consequence,
    \[\qq_p^\times /(\qq_p^\times)^2 \cong \begin{cases}
        (\zz/2\zz)^3, \quad & p  = 2 \\
        (\zz/2\zz)^2, \quad & p \neq 2
    \end{cases}\]
\end{corollary}
\begin{proof}
    
\end{proof}


\begin{theorem}[Hensel's lemma for polynomials]
    Let \((K,|\cdot|)\) be a complete non-archimedean normed field, and we assume that \(f \in \mathcal{O}_K[X]\) is a polynomial with a factorization over residue field \(k\)
    \[\bar{f} = \bar{g}\bar{h}\]
    with \(\bar{g},\bar{h}\) coprime in \(k[X]\), then there exists a lifting such that 
    \[f = gh\]
    with \(g = \bar{g} \mod \mathfrak{m}\), \(h = \bar{h} \mod \mathfrak{m}\) and \(\deg g = \deg \bar{g}\).
\end{theorem}

\begin{proof}
    
\end{proof}

\begin{corollary}[Hensel-Kurschak]
    Let \((K,|\cdot|)\) be a complete non-archimedean normed field, if \(f(X) = \sum_{i=0}^n a_iX^i \in K[X]\) is a irreducible polynomiakl, then 
    \[|a_i| \leq \max{|a_0|,|a_n|}\]
    for all \(i = 0,1,\dots,n\).
\end{corollary}

\begin{proof}
    
\end{proof}

\newpage

\subsection{Extension of absolute value
}

\begin{lemma} \label{equivalent-ultrametric}
    Let \(|\cdot|\) be an absolute value on the field \(K\), then the following conditions are equivalent:\\
    (1) \(|\cdot|\) is non-archimedean. \\
    (2) For any \(n \in \nn\), \(|n\cdot 1_K| \leq 1\).\\
    (3) For any \(x \in K\), if \(|x| \leq 1\) then \(|1+x| \leq 1\).\\
    (4) The closed unit ball in \(K\) is a subring.
\end{lemma}

\begin{proof}
    
\end{proof}

For a normed field \(K\), if we consider its field extension \(L/K\), then a natural question is the extension of the absolute value and its completion.

\begin{theorem}
    Let \((K,|\cdot|_K)\) be a complete non-archimedean normed field, and \(L/K\) is a finite extension of degree \(n\), then
    \[|y|_L : = |\mathrm{Norm}_{L/K}(y)|^{1/n}_K\]
    extends the absolute value \(|\cdot|_K\) such that \(L\) is a complete normed field.
\end{theorem}
\begin{proof}
    It is enough to show that \(\mathcal{O}_L\) is the integral closure of \(\mathcal{O}_K\) in \(L\), so it is a subring and we finish the proof by lemma \ref{equivalent-ultrametric}.


\end{proof}

\begin{remark}
    To understand such a construction, we assume that \(L/K\) is a Galois extension, or \(K \subset L \subset E\) a tower of extension with \(E/K\) galois. We assume \(|\cdot|_L\) is an extension of aboslute value, then for any \(\sigma \in \mathrm{Gal}(L/K)\), \(|\sigma(\cdot)|_L\) is also an extension. By the uniqueness of completion, we can conclude that \(|\cdot|_L = |\sigma(\cdot)|_L\), which allows us to calculate for any \(y \in L\)
    \begin{align*}
        |N_{L/K}(y)|_K = |N_{L/K}(y)|_L = |\prod_{\sigma} \sigma(y)|_L = \prod_\sigma |\sigma(y)|_L = \prod_\sigma |y|_L = |y|_L^{n}
    \end{align*}
\end{remark}

\begin{proposition}[Primitive element theorem]
    Let \(L/K\) be a finite extension of a complete non-archimedean normed field \(K\), assume that\\
    (1) \(\mathcal{O}_K\) is compact.\\
    (2) The extension of residue field \(k_L/k\) is finite and separable.\\
    Then  there exists an element \(\alpha \in L\) such that \(\mathcal{O}_L= \mathcal{O}_K[\alpha]\).
    
\end{proposition}

\begin{proof}
    
\end{proof}


